% !TeX root = ../../tesi.tex

\section{Realizzazione}

\subsection{L’agente conversazionale}

Descrivo il processo con cui l’agente costruisce una risposta: un ciclo in cui il modello invoca gli
strumenti per consultare il catalogo, finché non seleziona un componente da mostrare. Discuto come
ho ripartito il controllo del comportamento tra le istruzioni al modello e le regole fisse nel
codice, come ho generato i tre messaggi suggeriti e come ho scelto il modello e il livello di
ragionamento, confrontandone velocità e accuratezza. Un diagramma di sequenza segue il percorso di
un messaggio dall’utente alla risposta. Questa parte risponde agli obiettivi OB2 e OB3.

\subsection{Il \emph{widget} e l’integrazione con il tema del negozio}

Presento il \emph{widget}: i componenti (schede prodotto con selezione della variante, schede con le
risposte tratte dalle pagine del negozio, collegamenti alle collezioni), il supporto a più lingue e
valute e, in particolare, la sincronizzazione del carrello con il tema del negozio. Analizzo perché
ogni tema Shopify gestisce il carrello in modo diverso e come ho affrontato il problema con un
\emph{adapter} per ciascun tema e una soluzione di ripiego. Questa parte risponde agli obiettivi
OB1, OB4 e DE1.

\subsection{Metriche di utilizzo e analisi delle conversazioni}

Descrivo gli eventi che il \emph{widget} registra, il \emph{funnel} che la \emph{dashboard} presenta
al commerciante e l’analisi delle conversazioni con Jev, un modello decisionale di TypeSafe AI
uscito durante lo \emph{stage}: a differenza di un modello linguistico generativo, non produce testo
ma restituisce decisioni tipizzate con la loro probabilità. Da qui apro il confronto tra i due tipi
di modello nei compiti decisionali. Questa parte risponde all’obiettivo OB5.
